{\large\textbf{A Mechanical Model for Phase Transitions}}\footnote{Sitikantha Das (IIT Kharagpur) and Pramendra Ranjan Singh (Principal, Narayan College, J.P. University) were the principal authors of this problem. The contributions of the Academic Committee, Academic Development Group, and the International Board are gratefully acknowledged.}

A ring of radius $R$ has a bead of negligible size and mass $m$ threaded on
it. The ring is set rotating about its vertical diameter with angular velocity
$\omega$ as shown in Fig. 1. Alongside this, there is an opposing force on the
bead, $F_f$, which is proportional to the normal reaction $N$, and is given by
$F_f = fkN$ where $f$ is $1$ or $-1$. You are advised to employ polar
coordinates \{$r,\ \theta$\}. As far as possible express your answers in terms
of $\omega_c = \sqrt{g/R}$ where $g$ is the magnitude of the acceleration due
to gravity.

In a phase transition the free energy, $V(\mathcal{M})$ [mechanical equivalent
is potential energy] depends on the magnetization $\mathcal{M}$ as follows
\[
V(\mathcal{M}) = a(T)\mathcal{M}^2 + b(T)\mathcal{M}^4
\]
where $T\!\rightarrow$ temperature, $b(T) > 0$ and $a(T)$ changes sign with
temperature. We attempt to understand the phenomenon of phase transition using
the above mentioned model.

Note 1: For circular motion of radius $R$, the velocity in polar coordinates is
$\dot{\vec{r}} = R\dot{\theta}\hat{\theta}$ and the acceleration is
$\ddot{\vec{r}} = -R\dot{\theta}^2\hat{r} + R\ddot{\theta}\hat{\theta}$. Here
$\hat{r}$ and $\hat{\theta}$ are unit vectors in the radial and the tangential
directions respectively.

Note 2: The direction of the opposing force say $\vec{F}_f$ will be denoted by
$f$. Here $f$ = +1 if the bead is moving in the counter-clockwise direction (of
increasing) $\theta$ and $f$ = -1 if the bead is moving in the clockwise
direction $\theta$, e.g. $f$ = $sgn(\dot{\theta})$ where $sgn$ is +1 or -1
depending on whether its argument is positive or negative.

Note 3: You may find the expansions
\[
\begin{array}{rcl}
\sin(\theta) & = & \theta - \theta^3/6 + ..\\
\cos(\theta) & = & 1 - \theta^2/2 + \theta^4/24 + ..\\
(1+x)^n & = & 1 + nx + n(n-1)x^2/2 + n(n-1)(n-2)x^3/6 + ..
\end{array}
\]
(where $\theta$ is in radians and $|x| << 1$ ) useful for some parts of the
problem.

\begin{center}\includegraphics[width=0.5\linewidth]{figures/APhO_2022_Q2_fig1.png}\end{center}

\begin{center}Fig. 1: The bead on a rotating ring.\end{center}

In what follows we shall understand the dynamics of the bead in the frame of
the rotating ring and for angles in the range $-\pi/2 < \theta < \pi/2$. The
free body diagram of the bead is shown in Fig.2. Neglect all forces other than
the ones shown in the free body diagram.

\begin{center}\includegraphics[width=0.5\linewidth]{figures/APhO_2022_Q2_fig2.png}\end{center}

\begin{center}Fig. 2: The free body diagram.\end{center}

\textbf{A.1}\quad Write down the equations of motion for the radial $F_r$ and
the tangential $F_\theta$ components of the force on the bead. Assume $\theta$
increasing in the counter-clockwise direction. \hfill 0.5pt

\textbf{For the following parts B.1 to B.9 assume $k = 0$.}

\textbf{B.1}\quad State the relation between the equilibrium angle(s)
$\theta_0$ in terms of \{ $\omega,\ \omega_c$\}. \hfill 1.0pt

\textbf{B.2}\quad Qualitatively sketch $\theta_0$ ($y$-axis) as a function of
$\omega/\omega_c$ ($x$-axis). \hfill 0.5pt

\textbf{B.3}\quad Qualitatively sketch the magnitude of the normal reaction
force on the bead as a function of $\omega/\omega_c$ at stable equilibrium.
\hfill 0.5pt

\textbf{B.4}\quad We define the potential energy corresponding to the
tangential force $F_\theta$, namely
\[
F_\theta = -\frac{1}{R}\frac{d}{d\theta}V(\theta)
\]
with the zero of potential energy at $\theta = 0$. If $V(\theta)$ is expressed
as $P + Q\cos(\theta) + S\sin^2(\theta)$, obtain $P, Q, S$. \hfill 1.0pt

\textbf{B.5}\quad One can expand $V(\theta)$ for small $\theta$ and express it
as $V(\theta) = a(\omega)\,\theta^2{+}b(\omega)\,\theta^4$. Obtain the
coefficients $a(\omega)$ and $b(\omega)$. \hfill 1.0pt

\textbf{B.6}\quad Make representative plots of $V(\theta)$ versus $\theta$ for
values of $\omega/\omega_c$ just less than 1.0 (e.g. say 0.9) and
$\omega/\omega_c$ large, say 5.0. Note that only qualitative sketches and no
detailed calculations for the plots are required. \hfill 1.0pt

\textbf{B.7}\quad Landau theory of second order phase transitions can be used
to demarcate simple magnetic systems into two phases. For temperatures $T$
greater than the critical temperature $T_c$ the system is paramagnetic. For
$T < T_c$ the system is ferromagnetic and the magnetization $\mathcal{M}$ is
given by
\[
\mathcal{M}(T) = \mathcal{M}_0(1 - T/T_c)^{1/2} \quad T < T_c
\]
Let us denote the exponent $1/2$ by $\beta$. Compare this behaviour with the
bead problem discussed above. What are the analougues of $\mathcal{M}$, $T_c$,
$T/T_c$ in our case? What is the equivalent value of $\beta$ in our case?
\hfill 1.0pt

\textbf{B.8}\quad Determine the angular frequency of oscillation $\Omega_0$ of
the bead when it is disturbed from its \textit{``equilibrium''} position
$\theta_0$. Note that for small oscillations
\[
\Omega_0 = \frac{1}{R}\sqrt{\frac{V''(\theta_0)}{m}}
\]
\hfill 1.0pt

\textbf{B.9}\quad Qualitatively sketch $\Omega_0$ as a function of $\omega$.
\hfill 1.0pt

\textbf{For the following parts C.1 to C.2 $k \neq 0$.}

\textbf{C.1}\quad Take $f$ = 1 and express $k = \tan\alpha$. We may express
the condition for the equilibrium angle(s) $\theta_0$ as
\[
\left(\frac{\omega}{\omega_c}\right)^2 = \frac{\tan(x)}{\sin(y)}
\]
Obtain $x$ and $y$. \hfill 1.0pt

\textbf{C.2}\quad It is given that $f$ =1 and $k = 0.05$. Obtain the
equilibrium angles $\theta_0$, if any, for the following cases:\\
1.\ $\omega/\omega_c$ = 0.50\\
2.\ $\omega/\omega_c$ = 0.70 \hfill 0.5pt
